\chapter{The Morse Index Theorem}
\label{chap:dc-morse-index-theorem}

\section{The Index Theorem}

Let $\gamma:[0,a]\to M^n$ be a geodesic and let
\[
  \mathcal V=\mathcal V(0,a)=\{V:\text{piecewise differentiable along }\gamma:
  V(0)=V(a)=0\}.
\]

\begin{definition}[Index form]
  \label{def:dc-ch11-2-1}
  \uses{prop:dc-ch4-2-5}
  \dcref{ch11:2.1}
  The index form is the symmetric bilinear form
  \[
  I_a(V,W)=\int_0^a\bigl(\langle V',W'\rangle-
  \langle R(\gamma',V)\gamma',W\rangle\bigr)\,dt,
  \qquad V,W\in\mathcal V.
  \]
  Its index is the maximal dimension of a subspace on which $I_a(V,V)$ is
  negative definite. Its null space is
  \[
  \{V\in\mathcal V:I_a(V,W)=0\text{ for every }W\in\mathcal V\},
  \]
  and the nullity is its dimension. The form is degenerate when this dimension
  is positive. Symmetry follows from part (d) of \cref{prop:dc-ch4-2-5}.
\end{definition}

\begin{theorem}[Morse index theorem]
  \label{thm:dc-ch11-2-2}
  \dcref{ch11:2.2}
  The index of $I_a$ is finite and equals the number of points $\gamma(t)$,
  $0<t<a$, conjugate to $\gamma(0)$ along $\gamma$, with each point counted
  according to its multiplicity.
\end{theorem}
\begin{proof}
  \sketch
  For $t\in[0,a]$, let $i(t)$ be the index of the index form on
  $\gamma|_{[0,t]}$. By \cref{cor:dc-ch11-2-6}, \cref{lem:dc-ch11-2-7}, and
  \cref{lem:dc-ch11-2-8}, $i$ is zero near the origin, left-continuous, and its
  jump at $t$ equals the nullity of $I_t$. By \cref{cor:dc-ch11-2-4}, this nullity
  is the multiplicity of the conjugate point $\gamma(t)$; away from conjugate
  points it is zero. Thus $i(a)$ is the sum of the multiplicities of all conjugate
  points in $(0,a)$, proving \cref{thm:dc-ch11-2-2}.
\end{proof}

\begin{proposition}[Null space of the index form]
  \label{prop:dc-ch11-2-3}
  \dcref{ch11:2.3}
  An element $V\in\mathcal V$ lies in the null space of $I_a$ if and only if
  $V$ is a Jacobi field along $\gamma$.
\end{proposition}
\begin{proof}
  \sketch
  If $0=t_0<\cdots<t_k=a$ is a subdivision on which $V$ is differentiable, integration
  by parts gives
  \[
  I_a(V,W)=-\int_0^a\langle V''+R(\gamma',V)\gamma',W\rangle\,dt
  -\sum_{j=1}^{k-1}\left\langle V'(t_j^+)-V'(t_j^-),W(t_j)\right\rangle.\tag{1}
  \]
  A Jacobi field therefore annihilates $I_a$. Conversely, if $I_a(V,W)=0$ for all
  $W$, choose a differentiable $f$ positive away from the subdivision points and
  zero at them, and put $W=f(V''+R(\gamma',V)\gamma')$. Equation (1) yields
  \[
  \int_0^a f\,\lVert V''+R(\gamma',V)\gamma'\rVert^2\,dt=0,
  \]
  so $V$ satisfies the Jacobi equation on each open subinterval. A field $T$ with
  $T(t_j)=V'(t_j^+)-V'(t_j^-)$ then gives
  \[
  0=-\sum_{j=1}^{k-1}\lVert V'(t_j^+)-V'(t_j^-)\rVert^2,
  \]
  hence $V$ is $C^1$ and the Jacobi equation makes it smooth on all of $[0,a]$.
\end{proof}

\begin{corollary}[Degeneracy and conjugate endpoints]
  \label{cor:dc-ch11-2-4}
  \dcref{ch11:2.4}
  The form $I_a$ is degenerate exactly when $\gamma(a)$ is conjugate to
  $\gamma(0)$ along $\gamma$; then its nullity equals the multiplicity of
  $\gamma(a)$. Choose a normal subdivision
  \[
    0=t_0<t_1<\cdots<t_{k-1}<t_k=a
  \]
  for which every segment $\gamma|_{[t_{j-1},t_j]}$ lies in a totally normal
  neighborhood. Let $\mathcal V^-$ consist of fields that are Jacobi on each open
  subinterval and let $\mathcal V^+$ consist of fields vanishing at every interior
  subdivision point. The space $\mathcal V^-$ is finite-dimensional.
\end{corollary}

\begin{proposition}[Orthogonal splitting of the variation space]
  \label{prop:dc-ch11-2-5}
  \dcref{ch11:2.5}
  \mathcal V=\mathcal V^+\oplus\mathcal V^-,
  \qquad I_a(\mathcal V^+,\mathcal V^-)=0,
  \]
  and $I_a$ restricted to $\mathcal V^+$ is positive definite.
\end{proposition}
\begin{proof}
  \sketch
  For $V\in\mathcal V$, solve the Jacobi equation on each normal subinterval with
  the prescribed values $V(t_j)$; the resulting $W\in\mathcal V^-$ is unique and
  $V-W\in\mathcal V^+$. Integration by parts shows that for $X\in\mathcal V^-$ and
  $Y\in\mathcal V^+$ all boundary terms vanish, so $I_a(X,Y)=0$. Each normal
  segment minimizes energy, hence $I_a(Y,Y)\geq0$ for $Y\in\mathcal V^+$. If equality
  held for a nonzero $Y$, then for every $Z\in\mathcal V^+$ the inequality
\[
    0\leq I_a(Y+cZ,Y+cZ)=2cI_a(Y,Z)+c^2I_a(Z,Z)
  \]
  for all $c\in\mathbb R$ would force $I_a(Y,Z)=0$; orthogonality gives the same
  for $Z\in\mathcal V^-$. Thus $Y$ is a Jacobi field by \cref{prop:dc-ch11-2-3}, and
  its zeros at the subdivision points imply $Y=0$, a contradiction.
\end{proof}

\begin{corollary}[Finite-dimensional reduction]
  \label{cor:dc-ch11-2-6}
  \dcref{ch11:2.6}
  The index and nullity of $I_a$ equal those of its restriction to the finite-dimensional
  space $\mathcal V^-$. In particular, the index is finite.
\end{corollary}

For $t\in[0,a]$, write $\gamma_t=\gamma|_{[0,t]}$, let $I_t$ be its index form,
and let $i(t)$ denote the index of $I_t$. Choosing a normal subdivision with
$t\in(t_{j-1},t_j)$ identifies $\mathcal V^-(0,t)$ with
\[
  S_j=T_{\gamma(t_1)}M\oplus\cdots\oplus T_{\gamma(t_{j-1})}M.
\]
Thus the forms $I_t$ for $t\in(t_{j-1},t_j)$ act on a fixed finite-dimensional
space and vary continuously with $t$. The index is zero near $0$ and is
nondecreasing: extending a negative subspace on $[0,t]$ by zero to $[0,\bar t]$
preserves its form when $\bar t>t$.

\begin{lemma}[Left continuity of the index]
  \label{lem:dc-ch11-2-7}
  \dcref{ch11:2.7}
  If $\varepsilon>0$ is sufficiently small, then
  \[
    i(t-\varepsilon)=i(t).
  \]
  If $d$ denotes the nullity of $I_t$, then $d=0$ whenever $\gamma(t)$ is not
  conjugate to $\gamma(0)$.
\end{lemma}
\begin{proof}
  \sketch
  Monotonicity gives $i(t-\varepsilon)\leq i(t)$. On a subspace of dimension
  $i(t)$ where $I_t$ is negative definite, continuity of the family $I_s$ on the
  fixed space $S_j$ preserves negative definiteness for $s=t-\varepsilon$ close
  to $t$, giving the reverse inequality.
\end{proof}

\begin{lemma}[Index jump at a conjugate point]
  \label{lem:dc-ch11-2-8}
  \uses{thm:dc-ch11-2-2}
  \dcref{ch11:2.8}
  If $d$ is the nullity of $I_t$, then for sufficiently small $\varepsilon>0$,
  \[
    i(t+\varepsilon)=i(t)+d.
  \]
\end{lemma}
\begin{proof}
  \sketch
  Since $\dim S_j=n(j-1)$, $I_t$ is positive definite on a subspace of dimension
  $n(j-1)-i(t)-d$. Continuity preserves this subspace for $t+\varepsilon$, so
  $i(t+\varepsilon)\leq i(t)+d$. For the opposite inequality use the Index Lemma
  (\cref{lem:dc-ch10-2-2}): for a broken Jacobi field $V_{t_0}$ vanishing at
  $t_0\in(t_{j-1},t_j)$, extend it by zero on $[t_0,t_0+\varepsilon]$. The Index
  Lemma gives
  \[
    I_{t_0}(V_{t_0},V_{t_0})>
    I_{t_0+\varepsilon}(V_{t_0+\varepsilon},V_{t_0+\varepsilon}),
  \]
  with strictness because the extension is not Jacobi. Hence every vector with
  $I_t(V,V)\leq0$ becomes negative for $I_{t+\varepsilon}$; a negative subspace
  together with the null space of $I_t$ is therefore negative, yielding
  $i(t+\varepsilon)\geq i(t)+d$.
\end{proof}


\begin{corollary}[Jacobi's criterion]
  \label{cor:dc-ch11-2-9}
  \dcref{ch11:2.9}
  Let $\gamma:[0,a]\to M$ be a geodesic with $\gamma(a)$ not conjugate to
  $\gamma(0)$. Then $\gamma$ has no conjugate point in $(0,a)$ if and only if every
  proper variation has $E(s)\geq E(0)$ for all sufficiently small $|s|$. In
  particular, a minimizing geodesic has no conjugate point in its interior.
\end{corollary}

\begin{corollary}[Discreteness of conjugate points]
  \label{cor:dc-ch11-2-10}
  \dcref{ch11:2.10}
  The conjugate points along a geodesic form a discrete subset of the geodesic.
\end{corollary}

\begin{remark}[Finite-dimensional model of the path space]
  \label{rem:dc-ch11-2-12}
  \uses{rem:dc-ch9-2-7, cor:dc-ch11-2-6}
  \dcref{ch11:2.12}
  For the path space $\Omega_{p,q}$ and $c>0$, let $\Omega^c$ (respectively
  $\mathring{\Omega}{}^c$) consist of paths of energy at most $c$ (respectively
  strictly less than $c$). These paths lie in a compact set $S\subset M$. Choose
  $\delta>0$ so that points of $S$ at distance below $\delta$ have a unique
  minimizing geodesic, and subdivide $[0,1]$ by
  $|t_i-t_{i-1}|<\delta^2/c$. Let $B\subset\Omega^c$ be the broken geodesics
  whose pieces join consecutive subdivision points, and let $\mathring B$ be its
  intersection with $\mathring{\Omega}{}^c$. Since
  \[
    L^2(w_i)=(t_i-t_{i-1})E(w_i)<\delta^2,
  \]
  each piece is determined by its endpoints, and
  \[
    w\longmapsto(w(t_1),\ldots,w(t_{k-1}))
  \]
  identifies $\mathring B$ with an open subset of $M^{k-1}$. Its tangent vectors
  are broken Jacobi fields. There is a deformation retract
  $h_s:\mathring{\Omega}{}^c\to\mathring{\Omega}{}^c$ with image $\mathring B$.
  The restriction of the energy to $\mathring B$ has precisely the geodesics as
  critical points, and at a geodesic its Hessian has the same index and nullity as
  the index form by \cref{cor:dc-ch11-2-6}. Thus $\mathring B$ is a finite-dimensional
  model for the sublevel path space.
\end{remark}

\begin{remark}[Remark 2.11]
  \label{rem:dc-ch11-2-11}
  \dcref{ch11:2.11}
  The index theorem can be generalized to the case in which the points $\gamma(0)$
  and $\gamma(a)$ are replaced by submanifolds. For this see W. Ambrose [Am 2]. For
  an extremely general version of the index theorem, see S. Smale [Sm].
\end{remark}
